%% -*- coding: utf-8 -*-
\documentclass[12pt,pagesize,paper=192mm:108mm,landscape]{scrbook} 
%1920x1080 1280x720
\areaset[current]{192mm}{108mm}
\usepackage{calc}
\usepackage[T2A]{fontenc}
\usepackage[utf8]{inputenc}
\usepackage[english,russian]{babel}
\usepackage{microtype}
\usepackage{misccorr}
\usepackage{cmap}
%\usepackage[unicode=true]{hyperref}
\usepackage{graphicx}
\usepackage{amssymb}
\usepackage{amsmath}
%\usepackage{srcltx}
\usepackage{textcomp}
\usepackage{xspace}
%научные символы и смайлики \smiley \frownie
\usepackage{wasysym}
\usepackage{ccicons}
\begin{document}
\begin{titlepage}
  \vspace*{-0.5em}
  \begin{center}    
    \hspace*{3em}
    \begin{minipage}[t]{2em}
      \includegraphics[width=\textwidth]{../BINP-logo}
    \end{minipage}\hfill
    \begin{minipage}{0.12\linewidth}
    \includegraphics[width=\textwidth]{../NSU-logo}
    \end{minipage}
    \hfill
    \hspace*{5em}
    \smallskip

    Кафедра теоретической физики физического факультета НГУ

    \Large
    Профессор Грабовский А.\,В.

    \huge
    \textbf{Общая теория относительности}

    \Large
    Лекция № 4
    \vfill

    \normalsize
    \begin{minipage}{0.9\linewidth}
      Тензор и скаляр Риччи. Тензор Эйнштейна. Ковариантное сохранение
      тензора Эйнштейна. Принцип общей ковариантности. Ньютоновский
      предел. Выражение для $g_{00}$ из принципа эквивалентности в
      слабом статическом поле. Уравнения Эйнштейна. Космологическая
      постоянная. Тензор энергии-импульса, физический смысл его
      компонент: плотность энергии, компоненты плотности импульса,
      давление, плотность потока энергии, компоненты плотности потока
      импульса. Тензор энергии"=импульса пыли и идеальной
      жидкости. Ньютоновский предел уравнения Эйнштейна. Определение
      коэффициента пропорциональности в уравнении Эйнштейна.
     \end{minipage}
    \vfill

    \normalsize \ccbysa\hspace{0.5em}  Новосибирск 2022
  \end{center}
\end{titlepage}
\end{document}
